Because $A(l,E)$ and $B(l,E)$ have no singularities as functions of $l$, the only singularities (poles) of $S(l,E)$, and hence of the partial scattering amplitude $f(l,E)$, occur at zeros of $B(l,E)$. Poles of the scattering amplitude in the complex $l$-plane are termed \emph{Regge poles}. Their locations naturally depend on the value of the real parameter $E$. The functions $l=\alpha_i(E)$ that specify these locations are termed \emph{Regge trajectories}: as $E$ varies, the poles follow definite curves in the $l$-plane (below we shall omit the index $i$ that labels the poles).
